%=====================================================================
%  2. Related Work
%=====================================================================
\section{Related Work}
\label{sec:related}

\subsection{Intelligent Tutoring and Adaptive Learning}

Decades of research on intelligent tutoring systems (ITS) established that
adapting instruction to an explicit learner model produces reliable learning
gains \cite{vanlehn2011relative,kulik2016effectiveness}, and adaptive
hypermedia research formalized how content selection and sequencing should
depend on that model \cite{brusilovsky2001adaptive}. The ITS tradition is
methodologically strong on \emph{learner modeling} but weak on \emph{content
production}: instructional material is authored by humans and therefore cannot
scale to fast-moving technical domains or to long-tail learner profiles.

\subsection{Generative LLMs in Education}

The arrival of capable LLMs inverted this trade-off. Surveys document
widespread use for content creation, feedback, and tutoring
\cite{kasneci2023chatgpt,baidoo2023education}. Yet the dominant interaction
pattern remains a single model answering a single prompt, and the resulting
material tends to converge on one reading level and one rhetorical style
regardless of who asked. Two failure modes are repeatedly reported: factual
unreliability \cite{ji2023survey,huang2023survey} and the absence of any
mechanism that revises material in light of learner performance
\cite{doroudi2019instruction}. Retrieval-augmented generation
\cite{lewis2020rag,gao2023retrieval} reduces but does not eliminate the
former, because retrieval conditions the input without verifying the output.

\subsection{Multi-Agent LLM Systems}

Coordinating multiple LLM-driven agents has proved effective for complex
decomposition in software engineering \cite{hong2023metagpt,qian2023chatdev}
and general task solving \cite{wu2023autogen}. A parallel line of work shows
that adversarial interaction between models improves factuality and reasoning:
multi-agent debate \cite{du2023improving,liang2023encouraging}, self-refinement
loops \cite{madaan2023selfrefine,shinn2023reflexion}, and self-consistency
sampling \cite{wang2022selfconsistency}. Evaluation methodology has advanced
in step, with atomic factual precision metrics \cite{min2023factscore},
retrieval-aware scoring \cite{es2023ragas}, and model-based judging
\cite{zheng2023judging,chiang2024chatbotarena} now standard.

In education specifically, multi-agent work has concentrated on
\emph{simulation}: agents role-play teacher, student, or peer to generate
instructional dialogue \cite{zhang2024simulating}. These systems demonstrate
that agent orchestration can produce pedagogically plausible interaction, but
they are evaluated on interaction quality, not on the verifiable correctness
and learner-fit of the artifacts delivered.

\subsection{Positioning}

Table~\ref{tab:positioning} summarizes the gap. Prior work optimizes one or
two of the three properties that domain skill training requires
simultaneously: a persistent, multi-dimensional learner model; multi-modal
resource production; and pre-delivery verification. \sysname{} is designed to
satisfy all three, and---unlike debate methods evaluated only on reasoning
benchmarks---to measure the verification stage's contribution inside a
working educational pipeline. We are explicit about scope: the system generates
six resource modalities (document, mind map, graded quiz, extended reading,
runnable code, video script) over eight specialist agents plus debate, review,
and adaptation stages. The application layer additionally supports guided
hands-on instructions and illustrative figures, which are outside the
generation graph evaluated here.

\begin{table*}[!t]
\centering
\caption{Positioning of \sysname{} relative to representative prior work.
$\circ$ = partially addressed.}
\label{tab:positioning}
\begin{tabular}{lccc}
\toprule
\textbf{Work} & \textbf{Learner model} & \textbf{Multi-modal output} & \textbf{Pre-delivery verification} \\
\midrule
Classical ITS \cite{vanlehn2011relative,brusilovsky2001adaptive} & \checkmark & $\circ$ (human-authored) & $\circ$ \\
Single-LLM tutoring \cite{kasneci2023chatgpt,baidoo2023education} & $\circ$ & \checkmark & $\times$ \\
RAG-based tutoring \cite{lewis2020rag,gao2023retrieval} & $\circ$ & \checkmark & $\circ$ (retrieval only) \\
Multi-agent debate \cite{du2023improving,liang2023encouraging} & $\times$ & $\times$ & \checkmark \\
LLM agent classroom \cite{zhang2024simulating} & $\circ$ & $\circ$ & $\times$ \\
\midrule
\textbf{\sysname{} (this work)} & \checkmark & \checkmark & \checkmark \\
\bottomrule
\end{tabular}
\end{table*}
